\section{Introduction}

Environmental degradation is a pervasive outcome of the fundamental trade-off between economic productivity and the preservation of natural resources. This tension manifests primarily through two dominant, yet problematic, modes of resource interaction: \emph{extractivism} and \emph{predation}. Extractivism operates under the premise of sustainable limits but often drifts into overuse, while predation entails the outright disregard of ecological and legal constraints, leading to irreversible damage.

The conflict between short-term economic incentives and long-term ecological stability creates a persistent governance challenge, often resulting in suboptimal equilibria where individual rationality undermines collective welfare. Conventional economic models have historically treated the environment as an exogenous factor—a limitless input or a sink for waste—failing to integrate ecological feedback into the core of production theory and decision-making. This epistemological gap is a root cause of the systematic undervaluation of natural capital and the inefficacy of many environmental policies.

In response, this chapter proposes a theoretical reconceptualization grounded in the concept of \emph{economic commensalism}. This framework seeks to endogenize environmental variables within economic models, not as constraints or externalities, but as active, valuable inputs in the production function. By aligning private incentives with ecological carrying capacity through rational, market-compatible mechanisms, economic commensalism aims to transform preservation from a moral imperative or a regulatory burden into a strategically advantageous choice for economic agents.

The primary objective of this chapter is to lay the epistemological and conceptual foundations for this approach. It begins by critically examining the limitations of prevailing economic paradigms in addressing environmental degradation. It then introduces key ecological concepts—such as predation, mutualism, and commensalism—and adapts them for economic analysis. Finally, it constructs a heuristic model to demonstrate how incorporating the environment (\(E\)) as an endogenous variable can recalibrate incentives, leading to more sustainable and efficient outcomes. This theoretical groundwork is essential for the formal modeling and empirical validation pursued in subsequent chapters.

% fim %